\section{Conclusion}

This article presented an interpretable unsupervised learning workflow for identifying financial peer groups and changes in their membership among JSE Top 40 equities at three decision dates. UMAP dimensionality reduction was combined with average-linkage, Manhattan-distance hierarchical clustering, and the fitted groups were evaluated both in the reduced representation and in the scaled financial features.

The separation visible in the maps was much weaker in the financial features, so visual separation alone is insufficient evidence of strong financial peer structure. Of the 39 instruments present throughout, 20 changed aligned group at least once. A t-SNE comparator recovered the persistent groups and agreed on 28 of these 39 histories, which makes UMAP a defensible exploratory choice rather than a superior one. The clearest substantive findings are a continuing precious-metals peer core and a cross-sector group characterised by high market capitalisation and low JSE turnover, while other groups disperse; the maps, membership histories and financial profiles therefore need to be interpreted together.

A forward test showed that the clusters do not replace sector classifications, since sector peers co-moved more strongly in the following six months. Exploratory conditional analysis indicates that within sectors the clusters identify subgroups with stronger subsequent co-movement, whereas companies grouped with other sectors share financial characteristics rather than return behaviour. The workflow therefore supports exploratory screening and the refinement of sector peer groups, and its profiles, fidelity-qualified SHAP explanations and distance screen give equity analysts structured reasons to investigate unexpected relationships. Claims about structural anomalies or investment performance require further validation, beginning with the held-out evaluation of the June 2026 clusters under the now-fixed procedures.
